
\section{Del marco internacional al expediente concreto}
OACI proporciona el marco internacional y normas y métodos recomendados dentro de su sistema; su existencia no convierte cada documento en una orden ejecutable por un taller. En la Unión Europea, el Reglamento básico y sus actos de desarrollo distribuyen funciones y requisitos. Para un problema se identifica jurisdicción, producto, matrícula, operación y actividad antes de escoger la disposición.

Una autoridad puede aprobar o supervisar sin ser autora de todas las instrucciones técnicas. El titular de diseño emite información dentro de su función; una organización desarrolla procedimientos dentro de su aprobación. Distinguir emisor, aprobador, usuario y supervisor evita atribuir a un manual funciones jurídicas que no tiene.
\section{Obligación, medio de cumplimiento y orientación}
Un requisito vinculante expresa qué debe satisfacerse dentro de su ámbito. Un AMC ofrece una vía aceptable de demostrar cumplimiento del requisito asociado; un GM ayuda a interpretarlo. Los medios alternativos tienen procedimientos específicos. No toda frase de un compendio EAR tiene la misma fuerza por aparecer en el mismo PDF.

Una especificación de certificación se relaciona con una base y un procedimiento. Una nueva enmienda no sustituye automáticamente toda base anterior. Un SB procede del fabricante y puede resultar obligatorio por incorporación en otro requisito o dato vinculante; su título por sí solo no resuelve la cuestión.
\section{Distribuir responsabilidades sin perder trazabilidad}
Un expediente académico indica quién conserva programa, planifica, ejecuta, certifica y revisa. No asigna todo a «EASA» ni todo al mecánico. La contratación de una tarea no elimina por sí sola todas las responsabilidades del contratante. La estructura exacta depende de régimen, operación y organizaciones.

La hoja de control es una vista del sistema documental; no crea una aprobación. Si un registro es inconsistente se conserva evidencia y se activa una ruta de aclaración y corrección. Un cálculo interno correcto no sustituye la demostración de ejecución exigible.

\section{Convenio de Chicago, IATA y el antecedente JAA}
El Convenio de Chicago establece el marco internacional del que forman parte los estándares y métodos recomendados de OACI. Esos instrumentos buscan armonización entre Estados; su relación con la norma aplicable se analiza a través de los compromisos y sistemas pertinentes. No son una orden de mantenimiento de un componente ni deben confundirse con el manual del fabricante. La explicación institucional de OACI sobre elaboración de SARPs ayuda a identificar esa función \cite{ICAO-SARPS}.

IATA es una asociación de compañías aéreas. Su actividad sectorial, estándares y programas pueden ser relevantes para operadores, pero no la convierten en autoridad que expide licencias Parte 66 o aprueba una reparación en nombre de AESA. Se identifica la naturaleza del instrumento y, cuando se usa, su relación contractual, técnica o regulatoria concreta \cite{IATA}.

Las JAA son un antecedente del sistema europeo actual de armonización y sus JAR. EASA explica la transición desde una cooperación de autoridades nacionales hacia el marco común. Una referencia histórica JAR-145 en un libro antiguo no se usa como si sustituyera la Parte 145 vigente. Se conserva su utilidad histórica y se verifica el texto actual para cada afirmación \cite{EASA-JAA}.

\repaso{¿Una guía publicada en una web oficial es necesariamente legislación?}{La procedencia identifica la fuente, pero se distingue naturaleza y alcance. La guía se usa junto al requisito, su versión y sus condiciones.}
